Ahora procedemos a dar un pequeño vistazo a la teoria de \textit{espacios de Alexandroff}.
\begin{Def}
   Una topología $\tau$ sobre un conjunto $X$ se dice de \textbf{Alexandroff} si para cualquier colección $\left( U_i\right)_{i\in I}$ de elementos de $\tau$ se tiene $\bigcap_{i\in I}U_i \in\tau$. En este caso decimos que $(X,\tau)$ es un \textbf{espacio de Alexandroff}.
\end{Def}
De esta manera, un espacio topológico es de \textit{Alexandroff} si la intersección arbitraria de abiertos es un abierto. En el \Cref{Ejm:Topologias} vimos un espacio topológico en el que no se tiene esta propiedad. Ahora damos ejemplos de espacios topológicos de Alexandroff.
\begin{Ejm}
   \phantom{o}
   \begin{itemize}
      \item Cualquier espacio finito es de Alexandroff.
      \item Cualquier espacio con topología finita (en particular los espacios triviales) es de Alexandroff.
      \item Cualquier espacio discreto es de Alexandroff.
      \item $\mathbb{R}$ junto con la topología generada por la familia de intervalos $\left\lbrace [x,\infty)\mid x\in \mathbb{R}\right\rbrace$, llamada la topología de colas (a derecha) abiertas y cerradas, es un espacio de Alexandroff.
   \end{itemize}
\end{Ejm}
\begin{Def}
   Sean $X$ un espacio topológico y $x\in X$. Decimos que una vecindad $V$ de $x$ es una \textbf{vecindad mínima} de $x$ si para cualquier otra vecidad $U$ de $x$ se tiene $V\subseteq U$. Igualmente, decimos que una base $\mathcal{B}$ de $X$ es una \textbf{base mínima} si para cualquier base $\mathcal{B}'$ de $X$ se tiene $\mathcal{B}\subseteq\mathcal{B}'$.
\end{Def}
Una propiedad clave de los espacios de Alexandroff es que cada punto tiene una vecindad mínima, y cada espacio tiene una base mínima. Más precisamente:
\begin{Prop}
   Sea $X$ un espacio de Alexandroff. Para cada $x\in X$, la intersección de todas las vecindades de $x$ es una vecindad mínima de $x$, que denotamos por $U_x$. La colección $\mathcal{B}_{\operatorname{min}}=\left\lbrace U_x\mid x\in X\right\rbrace$ es una base mínima de $X$.
\end{Prop}
Las vecindades y bases mínimas de un espacio de Alexandroff nos dan una condición suficiente para que el espacio sea disconexo, como lo muestra la siguiente proposición, cuya prueba puede encontrarse en \cite{jafarian2013}.
\begin{Prop}\label{Prop:Conexo_Alexandroff}
   Sea $X$ un espacio de Alexandroff. Si existe $x\in X$ tal que:
   \begin{itemize}
      \item $U_x$ es minimal en $\mathcal{B}_{\operatorname{min}}$, es decir, para todo $y\in X$, si $U_y\subseteq U_x$ entonces $U_x=U_y$, y,
      \item $U_x$ es maximal en $\mathcal{B}_{\operatorname{min}}$, es decir, para todo $y\in X$, si $U_x\subseteq U_y$ entonces $U_x=U_y$,
   \end{itemize}
   entonces $U_x$ es un clopen de $X$ y por tanto $X$ es disconexo.
\end{Prop}
Contar con una subbase de un espacio topológico nos permite ``ahorraronos trabajo'' para conocer si el espacio es de Alexandroff, pues para esto basta determinar si la intersección arbitraria de abiertos subbásicos es abierta. Más precisamente:
\begin{Prop}\label{Prop:Alexandroff_con_subbase}
   Sean $X$ un espacio topológico y $\mathcal{S}$ una subbase de $X$. Se tiene que $X$ es un espacio de Alexandroff si, y solo si, la intersección arbitraria de elementos de $\mathcal{S}$ es un abierto de $X$.
\end{Prop}
\begin{proof}
   $(\Rightarrow)$ Se sigue del hecho de que los elementos de $\mathcal{S}$ son abiertos de $X$.\\
   $(\Leftarrow)$ Supongamos que la intersección arbitraria de elementos de $\mathcal{S}$ es un abierto de $X$. Sea $\left( U_i\right)_{i\in I}$ una colección cualquiera de abiertos de $X$ y llamemos $U:=\bigcap_{i\in I}U_i$. Usamos la \Cref{Prop:Caracterizacion_abiertos} para ver que $U \subseteqab X$. Sea $x\in U$ cualquiera. Sea $i\in I$ cualquiera. Tenemos que $x\in U_i$ y $U_i=\bigcup_{B\in A_i}B$ para algún $A_i\subseteq \mathcal{B}_{\mathcal{S}}$. De este modo $x\in B_i$ para algún $B_i\in A_i$. Como en particular $B_i\in \mathcal{B}_\mathcal{S}$ entonces existen $n_i\in \mathbb{Z}^+$ y $S_{i_j}\in \mathcal{S}$ para todo $j\in\left\lbrace 1,\dots ,n_i\right\rbrace$ tales que $B_i=\bigcap_{j=1}^{n_i}S_{i_j}$. Llamemos $V:=\bigcap_{i\in I}B_i=\bigcap_{i\in I}\left( \bigcap_{j=1}^{n_i}S_{i_j}\right)$. Como $V$ es una intersección de elementos de $\mathcal{S}$, se tiene $V\subseteqab X$. Como $x\in B_i$ para todo $i\in I$, entonces $x\in V$. Además, para todo $i\in I$, como $B_i\in A_i$, entonces $B_i\subseteq \bigcup_{B\in A_i}B=U_i$, luego $V=\bigcap_{i\in I}B_i\subseteq \bigcap_{i\in I}U_i=U$. Con esto, $U\subseteqab X$ y por tanto $X$ es un espacio de Alexandroff.
\end{proof}
La siguiente proposición nos habla sobre propiedades de separación en espacios de Alexandroff.
\begin{Prop}\label{Prop:Separacion_Alexandroff}
   Sea $X$ un espacio de Alexandroff.
   \begin{itemize}
      \item X es $T_0$ si, y solo si, $U_x=U_y$ implica $x=y$ para cualesquiera $x,y\in X$.
      \item X es $T_1$ si, y solo si, $X$ es discreto, lo cual es equivalente a que $U_x=\{x\}$ para todo $x\in X$.
   \end{itemize}
\end{Prop}
El siguiente es un hecho bastante importante y útil de los espacios de Alexandroff.
\begin{Prop}
   Sean $X$ y $Y$ espacios topológicos. Si $X$ es un espacio de Alexandroff y $X$ y $Y$ son homeomorfos, entonces $Y$ es también un espacio de Alexandroff. 
\end{Prop}
En otras palabras, ser de Alexandroff es una propiedad topológica.\\\\
Por último, enunciamos una proposición que será utilizada más adelante. Su prueba puede verse en \cite{jafarian2013}.
\begin{Prop}
   Sea $X$ un espacio de Alexandroff. Tomemos 
   $$
   \mathcal{B}:=\left\lbrace U_x\mid x\in X \text{ y }U_x\text{ es minimal en }\mathcal{B}_{\operatorname{min}}\right\rbrace.
   $$
   Se tiene lo siguiente:
   \begin{itemize}
      \item[1.] Si $A$ es un subconjunto denso minimal de $X$ entonces existe una función sobreyectiva $f:\mathcal{B}\to A$ tal que $f\left( U_x\right)\in U_x$ para todo $U_x\in \mathcal{B}$. En particular,
         $$
         A\subseteq\left\lbrace x\in X\mid U_x \text{ es minimal en }\mathcal{B}_{\operatorname{min}}\right\rbrace.
         $$
      \item[2.] Recíprocamente, si $f:\mathcal{B}\to X$ es una función tal que $f\left( U_x\right)\in U_x$ para todo $U_x\in \mathcal{B}$, entonces $f\left( \mathcal{B}\right)$ es un subconjunto denso minimal de $X$.
   \end{itemize}
   En particular, si $A$ y $A'$ son subconjuntos densos minimales de $X$ entonces $|A|=|A'|$.
\end{Prop}
Aclaramos que $A$ es un subconjunto denso minimal de $X$ si para cualquier otro subconjunto denso $B$ de $X$, la inclusión $B\subseteq A$ implica $B=A$.
